%Notes from Cliff Toubes, transcribed with ChatGPT
%\documentclass[12pt,letterpaper]{article}
%\usepackage{amsmath, amssymb, amsthm}

%%%
%% Required packages:
%%
\usepackage{etex}
\usepackage{amsmath}
\usepackage{amssymb}
\usepackage{amstext}
\usepackage{amsthm}
\usepackage{fancyhdr,fancybox}
\usepackage{mathrsfs} % need for \mathscr command
\usepackage{multicol}
\usepackage[normalem]{ulem}
%\usepackage{pictex,pictexplus}
% \usepackage{pictexwd}
\usepackage{rotating}
\usepackage{url}
\usepackage{xfrac}
\usepackage{booktabs}
\usepackage{color,soul}
\usepackage{comment}

\usepackage{hyperref}
\hypersetup{
    colorlinks=true,     % false: boxed links; true: colored links
    linkcolor=blue,      % color of internal links
    citecolor=blue,      % color of links to bibliography
    filecolor=blue,      % color of file links
    urlcolor=blue        % color of external links
}


\usepackage{tikz}
\usetikzlibrary{backgrounds,calc}

%%
%% Common formatting commands
%%
\setlength{\parindent}{0in}
\setlength{\textwidth}{7in}
\setlength{\evensidemargin}{-0.25in}
\setlength{\oddsidemargin}{-0.25in}
\setlength{\parskip}{.5\baselineskip}
\setlength{\topmargin}{-0.5in}
\setlength{\textheight}{9in}

%               Problem and Part
\newcounter{problemnumber}
\newcounter{partnumber}
\newcounter{subpartnumber}
\newcommand{\Problem}{%
  \refstepcounter{problemnumber}%
  \setcounter{partnumber}{0}%
  \setcounter{subpartnumber}{0}%
  \item[\makebox{\hfil\textbf{\theproblemnumber.}\hfil}]%
  }
\newcommand{\InProblem}[2][1.6]{%
  \refstepcounter{problemnumber}%
  \setcounter{partnumber}{0}%
  \setcounter{subpartnumber}{0}%
  \textbf{\theproblemnumber.}\ \ \parbox[t]{#1in}{#2}%
}
\newcommand{\InSmallProblem}[1]{\InProblem[1.05]{#1}}
\newcommand{\NoProblem}[1]{%
  \mbox{\hfil\textbf{#1}\hfil}%
  }
\newcommand{\UnProblem}{%
  \refstepcounter{problemnumber}%
  \setcounter{partnumber}{0}%
  \setcounter{subpartnumber}{0}%
  \mbox{\hfil\textbf{\theproblemnumber.}\hfil}%
  }
\newcommand{\InFlexProblem}[1]{%
  \refstepcounter{problemnumber}%
  \setcounter{partnumber}{0}%
  \setcounter{subpartnumber}{0}%
  \textbf{\theproblemnumber.}\ \ #1%
}
%%  Part:
\newcommand{\Part}{%
  \renewcommand{\thepartnumber}{(\alph{partnumber})}%
  \refstepcounter{partnumber}
  \setcounter{subpartnumber}{0}%
  \item[(\alph{partnumber})]%
}
\newcommand{\InPart}[2][1.75]{%
  \renewcommand{\thepartnumber}{(\alph{partnumber})}%
  \refstepcounter{partnumber}%
  \setcounter{subpartnumber}{0}%
  (\alph{partnumber})\ \ \parbox[t]{#1in}{#2}%
}
\newcommand{\UnPart}{%
  \refstepcounter{partnumber}%
  \renewcommand{\thepartnumber}{(\alph{partnumber})}%
  \setcounter{subpartnumber}{0}%
  (\alph{partnumber})%
}
\newcommand{\InLargePart}[1]{\InPart[3.05]{#1}}
\newcommand{\InFullPart}[1]{\InPart[6.2]{#1}}
\newcommand{\InSmallPart}[1]{\InPart[1.05]{#1}}
%% SubPart:
\newcommand{\SubPart}{%
  \refstepcounter{subpartnumber}%
  \item[(\roman{subpartnumber})]%
  }
\newcommand{\InSubPart}[2][2.25]{%
  \stepcounter{subpartnumber}%
  (\roman{subpartnumber})\ \ \parbox[t]{#1in}{#2}%
}
\newcommand{\InSmallSubPart}[1]{\InSubPart[1.5]{#1}}
\newcommand{\InTinySubPart}[1]{%
  \stepcounter{subpartnumber}%
  (\roman{subpartnumber})\ \ \makebox{#1}%
}
\newcommand{\InMySubPart}[2][2.25]{%
  \stepcounter{subpartnumber}%
  \parbox[t]{#1in}{\makebox[0.3in][l]{(\roman{subpartnumber})}\ \ {#2}}%
}
\newcommand{\TrueFalse}{\textbf{T}\ \ \textbf{F}\ \ }
\newcommand{\TFPart}[1]{% 
  \stepcounter{partnumber}%
  \setcounter{subpartnumber}{0}%
  \item[(\alph{partnumber})] \TrueFalse\parbox[t]{5.5in}{#1}}

\newcommand{\Points}[1]{(#1 points) \ }
\newcommand{\Point}[1]{(#1 point) \ }


%% For Answers:
\newcounter{answernumber}
\newcommand{\Answer}{%
  \refstepcounter{answernumber}%
  \setcounter{partnumber}{0}%
  \setcounter{subpartnumber}{0}%
  \item[\makebox{\hfil\textbf{\theanswernumber.}\hfil}]%
  }


% Setting up the headers for the first page
%\chead{} % will default to what is typed in the file.
\fancyhead[L]{\textsc{Math 22a}}
\fancyhead[R]{\textsc{Fall 2024}}
\fancyfoot[R]{
	\vspace*{\fill}\footnotesize\textcopyright{} \emph{Harvard Math 22a}	
}
\renewcommand{\headrulewidth}{0pt}

%\fancyhead[C]{}
%	\chead{}	
%	\lhead{}
%	\rhead{}
%	\cfoot{}


% Putting copyright on the footer of each page
%\fancypagestyle{secondpage}{
%	\fancyhead[C]{TESTing again} % change this to be blank
%		%\chead{}	
%	\fancyhead[L]{bears} % change this to be blank
%	\fancyhead[R]{dogs} % change this to be blank
%	\fancyfoot[R]{ %keeps the footer the same
%		\vspace*{\fill}\footnotesize\textcopyright{} \emph{Harvard Math 22a}	
%	}
%}
%\renewcommand{\headrulewidth}{0pt}
%\pagestyle{fancy}
\pagestyle{empty}


%\lhead{}
%\rhead{}
%\rfoot{\vspace*{\fill}\footnotesize\textcopyright{} \emph{Harvard Math 22a}}
%\renewcommand{\headrulewidth}{0pt}
%\pagestyle{fancy}




%%
%% Common shortcut commands:
%%
\newcommand{\ds}{\displaystyle}
\newcommand{\sgn}{\operatorname{sgn}}
\renewcommand{\Pr}{\operatorname{P}}
\newcommand{\CPr}[3][{}]{\Pr_{\text{#1}}\left(#2\,|\,#3\right)}

\newcommand{\C}{\mathbb{C}}
\newcommand{\R}{\mathbb{R}}
\newcommand{\Z}{\mathbb{Z}}
\newcommand{\N}{\mathbb{N}}
\newcommand{\Q}{\mathbb{Q}}
\newcommand{\D}{\mathbb{D}}

\newcommand{\rref}{\operatorname{rref}}
\newcommand{\im}{\operatorname{im}}
\newcommand{\rank}{\operatorname{rank}}
\newcommand{\diag}{\operatorname{diag}}
\newcommand{\Span}{\operatorname{span}}
%\renewcommand{\vec}[1]{\mathbf{#1}}
\newcommand{\charpoly}{\mathscr{P}}
\newcommand{\nicefrac}[2]{\sfrac{#1}{#2}} % using xfrac package
\newcommand{\topology}{\mathcal{T}}
\newcommand{\Inn}{\operatorname{Inn}}

\newcommand{\blankbox}[2]{\fbox{\rule{#1}{0in}\rule{0in}{#2}}}

\newenvironment{myindentpar}[1]%
 {\begin{list}{}%
         {\setlength{\leftmargin}{#1}
          \setlength{\rightmargin}{\leftmargin}}%
         \item[]%
 }
 {\end{list}}
\newtheorem{theorem}{Theorem}
\newtheorem*{NStheorem}{Theorem}
\newtheorem{cor}[theorem]{Corollary}
\newtheorem*{claim}{Claim}
\newtheorem{defn}{Definition}

\newenvironment{amatrix}[1]{%
  \left(\begin{array}{@{}*{#1}{c}|c@{}}
}{%
  \end{array}\right)
}

%--------grstep
% For denoting a Gauss' reduction step.
% Use as: \grstep{\rho_1+\rho_3} or \grstep[2\rho_5 \\ 3\rho_6]{\rho_1+\rho_3}
\newcommand{\grstep}[2][\relax]{%
   \ensuremath{\mathrel{
       {\mathop{\longrightarrow}\limits^{#2\mathstrut}_{
                                     \begin{subarray}{l} #1 \end{subarray}}}}}}
\newcommand{\swap}{\leftrightarrow}



\section{{Quadratic Forms and Optimization, $\vec\S$7.2, 7.3}}

\chead{\large\textbf{Quadratic Forms and Optimization, $\vec\S$7.2, 7.3}}

%\begin{document}
\thispagestyle{fancy}

Review from previously:

\begin{enumerate}
    \item Symmetric matrices: $A^T = A$.
    \item Fundamental fact about transpose: $u \cdot Av = (A^T u) \cdot v$.
    \item If $A$ is symmetric:
    \begin{enumerate}
        \item Two eigenvectors with distinct eigenvalues are orthogonal.
        \item All eigenvalues are real.
    \end{enumerate}
    \item Definition: An $n \times n$ matrix $A$ is said to be orthogonally diagonalizable if there is a diagonal matrix $D$ and orthogonal matrix $P$ with $P^T = P^{-1}$ such that $A = PDP^{-1}$.
    \item Theorem: A matrix is orthogonally diagonalizable if and only if it is symmetric.
    \begin{enumerate}
        \item Thus, $A$ has $n$ real eigenvalues counting multiplicity.
        \item The eigenvectors form a complete, orthogonal basis.
    \end{enumerate}
    \item Spectral theorem: Let $A$ be a symmetric matrix. Fix an orthonormal basis $\{v_k\}$ of eigenvectors. Then
    \[
    A = \lambda_1 v_1v_1^T + \cdots + \lambda_n v_nv_n^T.
    \]
\end{enumerate}

\textbf{Definition}:
A \textbf{quadratic form on $\mathbb{R}^n$} is a function (denoted by $Q$) whose value at $x$ can be computed using the rule $Q(x) = x^T A x = x \cdot A x$, where $A$ is a symmetric matrix.

\begin{enumerate}
    \Problem Which of the following functions are quadratic forms (and which aren't)? For those that are, write the corresponding symmetric matrix.
    \begin{enumerate}
        \item $Q(x) = x \cdot x$\vfill
        \item $Q(x) = -x_1^2 - x_2^2 + x_3^2 + x_4^2$\vfill
        \item $Q(x) = x_1x_2 + x_3x_4$\vfill
        \item $Q(x) = x_1 - x_2x_3$\vfill
        \item $Q(x) = \cos(x_1^2 - x_2^2)$\vfill
    \end{enumerate}

\newpage

    \Problem Consider the quadratic form $Q(x) = x_1^2 - x_2^2$. Compute the quadratic form $y \mapsto Q(Py)$ where $P$ is the matrix below:
    \begin{enumerate}
        \item $P = \begin{pmatrix} 2 & 0 \\ 0 & \frac{1}{2} \end{pmatrix}$\vfill
        \item $P = \frac{1}{\sqrt{5}} \begin{pmatrix} 2 & 1 \\ 1 & 2 \end{pmatrix}$\vfill
    \end{enumerate}

    \Problem Diagonalize the following quadratic forms (find the diagonal matrix $D$ and the orthogonal change of variable matrix $P$):
    \begin{enumerate}
        \item $Q(x) = x_1^2 - 2x_1x_2 + x_2^2$\vfill
        \item $Q(x) = 2x_1^2 - 6x_1x_2 + 2x_2^2$\vfill
        \item $Q(x) = x_1x_2 + x_3x_4$\vfill
    \end{enumerate}

\newpage
    \Problem For each of the quadratic forms below, what are the minimum and maximum values on the set where $\|x\| = 1$? And at what points are they achieved? And what are the saddle points and the corresponding value for $Q$?
    \begin{enumerate}
        \item $Q(x) = x_1^2 - 4x_2^2$\vfill
        \item $Q(x) = 7x_1^2 + 4x_2^2 + 3x_3^2$\vfill
        \item $Q(x) = -7x_1^2 + 4x_2^2 + 3x_3^2$\vfill
        \item $Q(x) = -7x_1^2 - 4x_2^2 + 3x_3^2$\vfill
        \item $Q(x) = -7x_1^2 - 4x_2^2 - 3x_3^2$\vfill
    \end{enumerate}

\newpage

    \Problem Find the minimum and maximum values of the following quadratic forms on the set where $\|x\| = 1$; and give their coordinates. Also, find all saddle points with their coordinates.
    \begin{enumerate}
        \item $Q(x) = 2x_1x_2 + 2x_2x_3 + 2x_3x_1$\vfill
        \item $Q(x) = 2x_1x_2 - 2x_2x_3$\vfill
        \item $Q(x) = 4x_1x_2 - x_1^2 - x_2^2$\vfill
    \end{enumerate}
\end{enumerate}





\begin{comment}

%TODO
% -Format the below and check the answers

\newpage
\chead{\Large\textbf{Quadratic Forms and Optimization, - Answers / Solutions}}%
\lhead{}
\rhead{}
\thispagestyle{fancy}


\begin{enumerate}
    \item 
    \begin{enumerate}
        \item Yes, $A = I_n$.
        \item Yes, $A = \begin{pmatrix} -1 & 0 & 0 & 0 \\ 0 & -1 & 0 & 0 \\ 0 & 0 & 1 & 0 \\ 0 & 0 & 0 & 1 \end{pmatrix}$.
        \item Yes, $A = \begin{pmatrix} 0 & 1 & 0 & 0 \\ 1 & 0 & 0 & 0 \\ 0 & 0 & 0 & 1 \\ 0 & 0 & 1 & 0 \end{pmatrix}$.
        \item No.
        \item No.
    \end{enumerate}

    \item 
    \begin{enumerate}
        \item $Q(Py) = 4y_1^2 - \frac{1}{4} y_2^2$.
        \item $Q(Py) = y_1^2 - y_2^2$.
    \end{enumerate}

    \item 
    \begin{enumerate}
        \item $D = \begin{pmatrix} 2 & 0 \\ 0 & 0 \end{pmatrix}$, $P = \frac{1}{\sqrt{2}} \begin{pmatrix} 1 & 1 \\ -1 & 1 \end{pmatrix}$.
        \item $D = \begin{pmatrix} 5 & 0 \\ 0 & -1 \end{pmatrix}$, $P = \frac{1}{\sqrt{2}} \begin{pmatrix} 1 & 1 \\ -1 & 1 \end{pmatrix}$.
        \item $D = \frac{1}{2} \begin{pmatrix} 1 & 0 & 0 & 0 \\ 0 & -1 & 0 & 0 \\ 0 & 0 & 1 & 0 \\ 0 & 0 & 0 & -1 \end{pmatrix}$, $P = \frac{1}{\sqrt{2}} \begin{pmatrix} 1 & 1 & 0 & 0 \\ 1 & -1 & 0 & 0 \\ 0 & 0 & 1 & 1 \\ 0 & 0 & 1 & -1 \end{pmatrix}$.
    \end{enumerate}

    \item 
    \begin{enumerate}
        \item Maximum value = 1 where $x_2 = 0$, minimum = -4 where $x_1 = 0$.
        \item Maximum value = 7 where $x_2 = x_3 = 0$, minimum value = 3 where $x_1 = x_2 = 0$, saddle points where $x_1 = x_3 = 0$ with value 4.
        \item Maximum value = 4 where $x_1 = x_3 = 0$, minimum value = -7 where $x_2 = x_3 = 0$, saddle points where $x_1 = x_2 = 0$ with value 3.
        \item Maximum value = 3 where $x_1 = x_2 = 0$, minimum value = -7 where $x_2 = x_3 = 0$, saddle points where $x_1 = x_2 = 0$ with value -4.
        \item Maximum value = -3 where $x_1 = x_2 = 0$, minimum value = -7 where $x_2 = x_3 = 0$, saddle points where $x_1 = x_3 = 0$ with value -4.
    \end{enumerate}

    \item 
    \begin{enumerate}
        \item The matrix for $A$ is $\begin{pmatrix} 0 & 1 & 1 \\ 1 & 0 & 1 \\ 1 & 1 & 0 \end{pmatrix}$ which has eigenvalues $\lambda = 2$ with eigenvector $\begin{pmatrix} 1 \\ 1 \\ 1 \end{pmatrix}$, and $\lambda = -1$ with a 2-dimensional eigenspace spanned by $\begin{pmatrix} 1 \\ 0 \\ -1 \end{pmatrix}$ and $\begin{pmatrix} 0 \\ 1 \\ -1 \end{pmatrix}$. Thus, the maximum of $Q$ on the sphere where $\|x\| = 1$ is 2 at the points $\pm \frac{1}{\sqrt{3}} \begin{pmatrix} 1 \\ 1 \\ 1 \end{pmatrix}$, and the minimum is -1 on the circle of radius 1 in the plane where $x + y + z = 0$.
        \item The matrix for $A$ is $\begin{pmatrix} 0 & 1 & 0 \\ 1 & 0 & -1 \\ 0 & -1 & 0 \end{pmatrix}$ which has eigenvalues $\lambda = 0$ and $\lambda = \pm \sqrt{2}$. The maximum of $Q$ is therefore $\sqrt{2}$ and that is taken on at the points $\pm \frac{1}{2} \begin{pmatrix} 1 \\ \sqrt{2} \\ -1 \end{pmatrix}$, the minimum is $-\sqrt{2}$ which is taken on at $\pm \frac{1}{2} \begin{pmatrix} 1 \\ -\sqrt{2} \\ -1 \end{pmatrix}$, and then there are saddle points at $\pm \frac{1}{\sqrt{2}} \begin{pmatrix} 1 \\ 0 \\ 1 \end{pmatrix}$ where $Q$ is zero.
        \item The matrix $A$ is $\begin{pmatrix} -1 & 2 \\ 2 & -1 \end{pmatrix}$ which has eigenvalues $\lambda = 1$ and $\lambda = -3$. Thus, the maximum of $Q$ is 1 which occurs at the points $\pm \frac{1}{\sqrt{2}} \begin{pmatrix} 1 \\ 1 \end{pmatrix}$, and the minimum of $Q$ is -3 which occurs at $\pm \frac{1}{\sqrt{2}} \begin{pmatrix} 1 \\ -1 \end{pmatrix}$.
    \end{enumerate}
\end{enumerate}

\end{comment}

%\end{document}
